\documentclass[12pt,oneside]{article}
\usepackage[T1]{fontenc}
\usepackage{amsmath,amssymb,array,longtable}
\usepackage[english]{babel}
\usepackage[bookmarks=false,pdfauthor={Tomi Setsu},pdftitle={FunC v0.5.0 to Fift Lowering Guide}]{hyperref}
\usepackage{fancyhdr}

\setlength{\headheight}{15.2pt}
\pagestyle{fancy}
\fancyhf{}
\fancyfoot[C]{\thepage}
\renewcommand{\headrulewidth}{0.5pt}
\fancyhead[C]{FunC v0.5.0 to Fift Lowering Guide}

\newcolumntype{L}[1]{>{\raggedright\arraybackslash}p{#1}}

\title{FunC v0.5.0 $\rightarrow$ Fift\\Feature-by-Feature Lowering Guide}
\author{Tomi Setsu}

\begin{document}
\maketitle

\begin{abstract}
This document complements \verb|doc/func_v0.5.0.tex| by explaining how the
proposed FunC \verb|0.5.0| surface can be lowered into Fift/TVM code without
introducing a hidden runtime or a second virtual machine. The goal is to make
the lowering model explicit enough that language design can be judged against
the real constraints of TVM/TOS.

The central claim is that FunC \verb|0.5.0| should remain a \emph{front-end
heavy, runtime-light} language. Most new features are either:

\begin{enumerate}
\item compile-time only and erased before code generation;
\item mapped to fixed TVM stack or cell representations;
\item implemented by generated wrappers around existing TVM conventions.
\end{enumerate}

This is a design-lowering specification, not a statement that the current
compiler already implements all of the described transformations.
\end{abstract}

\section*{Introduction}
The proposed FunC \verb|0.5.0| surface is intentionally more expressive than
FunC \verb|0.4.6|: it adds named structs, enums, message declarations,
\verb|Option|/\verb|Result|, explicit mutability, explicit effect sets, traits,
generics, typed contexts, and typed cell APIs. The obvious concern is whether
these features remain faithful to the low-level execution model.

They can, provided that the compiler obeys three discipline rules:

\begin{enumerate}
\item no hidden heap or object runtime;
\item no dynamic dispatch requirement;
\item no source feature whose meaning cannot be reduced to ordinary TVM stack,
cell, control-flow, and action-phase operations.
\end{enumerate}

This document therefore answers a practical question: \emph{for each major
FunC 0.5.0 feature, what is the canonical lowering path to Fift/TVM?}

\clearpage
\tableofcontents

\clearpage
\section{Scope and Backend Model}
\subsection{What ``Lowering to Fift'' Means}
In this document, ``lowering to Fift'' means lowering to the same TVM execution
model already targeted by current FunC and Fift macro assembly. The backend may
internally pass through an IR, but the observable result should still be a
selector-based TVM program expressible in Fift assembly.

In particular:

\begin{itemize}
\item ordinary functions still become TVM routines callable through selector
indices and the existing \verb|CALLDICT|/\verb|JMPDICT| model;
\item cell parsing still becomes \verb|CTOS|, \verb|LDU|/\verb|LDUX|,
\verb|LDI|/\verb|LDIX|, \verb|LDREF|, and related instructions;
\item cell building still becomes \verb|NEWC|, \verb|STU|/\verb|STUX|,
\verb|STI|/\verb|STIX|, \verb|STREF|, and \verb|ENDC|;
\item outbound messaging still ends in \verb|SENDRAWMSG|;
\item persistent state still goes through the existing data-register
conventions and helpers such as \verb|get_data|/\verb|set_data|.
\end{itemize}

\subsection{Three Representation Layers}
Every source construct must be understood at three different layers:

\begin{enumerate}
\item \textbf{source representation}: what the author writes;
\item \textbf{runtime representation}: what exists on the TVM stack while code
executes;
\item \textbf{wire/storage representation}: how values are serialized into
cells for messages, persistent state, dictionaries, or ABI boundaries.
\end{enumerate}

Many confusions disappear once these layers are kept distinct. For example,
\verb|Option<T>| need not have the same representation on the stack and on the
wire, as long as the compiler-generated codecs are deterministic.

\subsection{Canonical Lowering vs Optimization}
This document specifies a \emph{canonical} lowering. Optimizations are still
allowed afterwards:

\begin{itemize}
\item tuples may be scalarized into separate locals;
\item enum tests may become jump tables or branch trees;
\item \verb|Option<T>| may be unboxed locally when the optimization is not
observable;
\item small helper functions may be inlined;
\item dead fields or dead temporary tuples may be removed.
\end{itemize}

The important invariant is that optimization may improve code generation, but
must not change the source-level meaning or the externally visible wire format.

\section{Lowering Principles}
\subsection{Compile-Time Features Must Erase}
Several proposed \verb|0.5.0| features exist to improve readability and static
safety, not to create a runtime subsystem. In particular:

\begin{itemize}
\item \verb|mod|, \verb|use|, and \verb|pub| are namespace features only;
\item effect sets are checked statically and erased;
\item mutable-place overlap checks are static only;
\item traits and generics use static dispatch and monomorphization;
\item field names and variant names are debug/source concepts, not runtime
symbols.
\end{itemize}

\subsection{Runtime Features Must Map to Existing TVM Values}
When a feature survives past type checking, it must lower to existing TVM value
kinds: integers, cells, slices, builders, tuples, continuations, or null.
Nothing in the proposal requires a new runtime type.

\subsection{Protocol-Aware Wrappers Are Allowed}
Some source forms correspond directly to blockchain protocol concepts rather
than to ordinary library code. These may use generated wrappers:

\begin{itemize}
\item \verb|entry internal|, \verb|entry external|, \verb|entry bounce|;
\item typed message contexts such as \verb|InternalContext<T>|;
\item \verb|message| declarations with bound opcodes;
\item \verb|BodyLayout::Auto| and typed \verb|send(... )|.
\end{itemize}

This is acceptable because these features do not invent a second semantic
world; they merely expose real TVM/TOS conventions in a safer surface form.

\subsection{Static Dispatch Is Mandatory}
To keep the backend simple and predictable:

\begin{itemize}
\item trait calls must resolve at compile time;
\item generic functions must be monomorphized;
\item there are no trait objects, no vtables, and no hidden dictionary-passing
runtime.
\end{itemize}

\section{Canonical Runtime Representations}
\subsection{Scalars and Domain Types}
\begin{longtable}{|L{0.24\textwidth}|L{0.30\textwidth}|L{0.32\textwidth}|}
\hline
\textbf{Source type} & \textbf{Canonical TVM stack representation} & \textbf{Wire/storage representation} \\ \hline
\endfirsthead
\hline
\textbf{Source type} & \textbf{Canonical TVM stack representation} & \textbf{Wire/storage representation} \\ \hline
\endhead
\texttt{iN}/\texttt{uN} &
ordinary TVM integer &
fixed-width or variable-width integer encoding, depending on codec \\ \hline
\texttt{bool} &
ordinary TVM integer, normalized to \texttt{0} for false and \texttt{-1} for true &
usually a 1-bit or fixed small-integer encoding chosen by the relevant codec \\ \hline
\texttt{Coins} &
ordinary TVM integer with source-level domain meaning &
\texttt{Grams}/\texttt{VarUInteger} encoding through the standard amount codecs \\ \hline
\texttt{Address} &
ordinary TVM \texttt{Slice} carrying the invariant ``valid serialized address'' &
serialized message address form \\ \hline
\texttt{RawCell} &
ordinary TVM \texttt{Cell} &
ordinary cell reference \\ \hline
\texttt{Cell<T>} &
ordinary TVM \texttt{Cell}; \texttt{T} is erased after type checking &
same raw cell bits and refs, but interpreted through \texttt{T}'s codec \\ \hline
\texttt{RawSlice} / \texttt{Slice} &
ordinary TVM \texttt{Slice} &
not directly stored; arises from \texttt{CTOS} or message/body parsing \\ \hline
\texttt{Builder} &
ordinary TVM \texttt{Builder} &
not directly stored; finalized with \texttt{ENDC} \\ \hline
\texttt{Bytes} &
byte-aligned \texttt{Slice} or stdlib wrapper over a slice/cell pair &
byte-string codec; exact layout must be fixed by the stdlib/ABI \\ \hline
\texttt{()} &
zero stack values at function boundaries; null/empty payload when embedded &
zero-width payload or omitted payload by codec rule \\ \hline
\end{longtable}

\subsection{Aggregates and Sum Types}
\begin{longtable}{|L{0.24\textwidth}|L{0.30\textwidth}|L{0.32\textwidth}|}
\hline
\textbf{Source construct} & \textbf{Canonical TVM stack representation} & \textbf{Wire/storage representation} \\ \hline
\endfirsthead
\hline
\textbf{Source construct} & \textbf{Canonical TVM stack representation} & \textbf{Wire/storage representation} \\ \hline
\endhead
\texttt{(T1, ..., Tn)} &
TVM tuple or flattened multi-value temporary &
codec-dependent tuple/product layout \\ \hline
\texttt{struct S \{...\}} &
logical tuple in field order; backend may scalarize locals &
derived product codec in field order \\ \hline
\texttt{enum E \{...\}} &
tagged tuple \texttt{(tag, payload)}; payload is a tuple or null &
tag plus payload, with tag width fixed by the derived codec \\ \hline
\texttt{Option<T>} / \texttt{T?} &
tagged tuple like any two-variant enum; optimizer may unbox locally &
\texttt{Maybe}-style tag plus payload \\ \hline
\texttt{Result<T, E>} &
tagged tuple \texttt{(0, ok)} or \texttt{(1, err)} canonically &
tag plus payload if explicitly serialized \\ \hline
\texttt{message M(op = ...)\{...\}} &
same as a named struct at runtime, plus compile-time opcode metadata &
opcode prefix followed by field encoding \\ \hline
\end{longtable}

\section{Feature-by-Feature Lowering Table}
\begingroup
\setlength{\tabcolsep}{4pt}
\footnotesize
\begin{longtable}{|L{0.15\textwidth}|L{0.24\textwidth}|L{0.24\textwidth}|L{0.15\textwidth}|}
\hline
\textbf{Feature} & \textbf{Frontend handling} & \textbf{Canonical Fift/TVM lowering} & \textbf{Notes} \\ \hline
\endfirsthead
\hline
\textbf{Feature} & \textbf{Frontend handling} & \textbf{Canonical Fift/TVM lowering} & \textbf{Notes} \\ \hline
\endhead
\texttt{mod}, \texttt{use}, \texttt{pub} &
resolve names, visibility, and imports at compile time &
no runtime code; generate mangled function/type symbols only &
pure erasure \\ \hline
\texttt{const}, type alias &
evaluate or resolve during semantic analysis &
constant immediates or no code at all &
pure erasure except constant uses \\ \hline
\texttt{let} &
create immutable local binding &
ordinary local slot / stack temporary &
no runtime mutability bit \\ \hline
\texttt{let mut} &
mark local as assignable &
same local slot shape as \texttt{let} &
\texttt{mut} is checker-only \\ \hline
assignment to locals/fields &
verify place mutability and field existence &
stack shuffles plus tuple rebuild or scalar local update &
backend may scalarize struct fields \\ \hline
\texttt{bool} &
type-check boolean-only operators and conditions &
emit ordinary integer-producing comparisons; normalize to \texttt{0/-1} &
no new VM type \\ \hline
\texttt{Coins} &
track domain meaning distinctly from raw ints &
ordinary TVM integer in computation; amount codec on serialization &
zero-cost in arithmetic paths \\ \hline
\texttt{Address} &
type-check address-only APIs &
ordinary slice value plus address-validity invariant &
codec uses message-address helpers \\ \hline
\texttt{Cell<T>} &
enforce typed cell APIs and erase \texttt{T} after checking &
ordinary cell value; codecs inserted at load/store boundaries &
no runtime generic metadata \\ \hline
\texttt{struct} &
resolve field offsets and patterns &
tuple-like product value or flattened locals &
field names erased after lowering \\ \hline
\texttt{enum} &
assign canonical tags and verify exhaustiveness &
tagged tuple \texttt{(tag, payload)}; branch on tag &
payload may be null for zero-field variants \\ \hline
\texttt{message(op = X)} &
derive codec and bind opcode metadata &
generated decoder checks prefix, then parses fields; encoder prepends opcode &
protocol-aware wrapper justified \\ \hline
\texttt{Option<T>} / \texttt{T?} &
desugar \texttt{T?} to \texttt{Option<T>} &
two-variant tagged tuple canonically &
optimizer may unbox locally \\ \hline
\texttt{Result<T, E>} &
type-check recoverable error paths &
two-variant tagged tuple canonically &
drives \texttt{?} lowering \\ \hline
\texttt{trait} / \texttt{impl} &
resolve methods and trait obligations statically &
no runtime object; selected concrete function symbol only &
no vtables \\ \hline
generics &
instantiate per concrete type vector &
monomorphized concrete functions/types, then ordinary lowering &
same discipline as template expansion \\ \hline
\texttt{mut self} methods &
desugar method call into explicit state threading &
internal helper returns updated receiver plus result, then uses existing TVM ops &
surface replacement for \texttt{\~{}method} \\ \hline
\texttt{mut} arguments &
verify mutable-place rules &
same calling convention as explicit state threading or by-value update &
overlap checks are static only \\ \hline
\texttt{if let} &
desugar to match over one scrutinee &
tag/type test plus destructuring branch &
for \texttt{Option}, often just tag test \\ \hline
\texttt{while let} &
desugar to loop plus match/test &
loop header, test, destructure, exit branch &
no special runtime form \\ \hline
\texttt{match} &
check exhaustiveness, bind arm patterns &
discriminant test tree, jump table, or direct destructuring &
message opcodes may use preload fast path \\ \hline
\texttt{?} on \texttt{Result} &
check enclosing residual type &
branch on tag; on error construct early return and jump to epilogue &
zero hidden exceptions \\ \hline
\texttt{effect(...)} &
check allowed operations transitively &
erased before code generation &
may survive only as debug metadata/comments \\ \hline
\texttt{extern "tvm"} &
bind a source signature to a raw intrinsic &
emit direct intrinsic/opcode sequence or imported asm helper &
surface replacement for parser magic \\ \hline
\texttt{entry internal} &
generate ABI wrapper around typed user function &
existing internal-entry ABI plus message/context parsing wrapper &
still one real TVM entrypoint \\ \hline
\texttt{entry external} &
generate ABI wrapper around typed user function &
existing external-entry ABI plus parse/signature wrapper &
\texttt{accept\_message} remains explicit unless specified otherwise \\ \hline
\texttt{entry bounce} &
generate wrapper for bounced-message path &
existing bounced-message convention plus bounce-body decoding wrapper &
no new protocol primitive \\ \hline
\texttt{get fn} &
mark as getter export and read-only surface &
ordinary getter selector / method-id export &
same getter ABI family as today \\ \hline
\texttt{InternalContext<T>} and friends &
type-check context field access and requested body type &
compiler-generated tuple/view assembled from real VM inputs &
context object itself is not heap state \\ \hline
\texttt{to\_cell}/\texttt{from\_cell} &
select \texttt{CellEncode}/\texttt{CellDecode} implementation &
\texttt{NEWC} + stores + \texttt{ENDC}; or \texttt{CTOS} + loads + \texttt{ENDS} &
typed facade over raw cell ops \\ \hline
\texttt{load\_any}/\texttt{store\_any} &
resolve concrete codec statically &
call monomorphized codec helper, which emits raw load/store instructions &
no reflective runtime \\ \hline
\texttt{send<T>} &
resolve body codec and send options &
encode body, build message envelope, choose inline/ref layout, emit \texttt{SENDRAWMSG} &
typed wrapper over action-cell construction \\ \hline
\texttt{BodyLayout::Auto} &
mark layout policy in source &
generated size-aware branch chooses inline body or ref body deterministically &
must be fully deterministic \\ \hline
\texttt{BounceMode} &
select bounce policy at source level &
maps to message-header bits and chosen bounced-body convention &
no dedicated VM primitive \\ \hline
\texttt{for item in items} &
desugar through a statically resolved iterator protocol &
hidden iterator state plus \texttt{loop} and \texttt{next()} match &
initial implementation should restrict iterable kinds \\ \hline
\end{longtable}
\endgroup

\section{Detailed Lowering Notes}
\subsection{Functions Still Lower to Selector-Based TVM Routines}
After desugaring, ordinary functions should still lower to the same broad model
used today:

\begin{enumerate}
\item each concrete function body is assigned a selector;
\item the program stores its selector dispatcher in \verb|c3|;
\item intra-program calls become \verb|CALLDICT n| or equivalent selector
calls;
\item tail calls become \verb|JMPDICT n| or equivalent jumps.
\end{enumerate}

The proposed language features therefore sit \emph{above} the current function
family model; they do not replace it.

\subsection{Explicit Mutability Lowers to Explicit State Threading}
A method with \verb|mut self| should not require a new calling convention at
the VM level. The canonical lowering is:

\begin{enumerate}
\item rewrite the surface method call into an internal helper call;
\item make the helper return both the updated receiver and the value result;
\item rebind the caller-side mutable place to the updated receiver;
\item keep using the existing raw TVM instructions underneath.
\end{enumerate}

For example:

\begin{verbatim}
let mut cs = body;
let op = cs.load_uint(32)?;
\end{verbatim}

canonically lowers in two conceptual steps:

\begin{verbatim}
let (__cs1, __r1) = __Slice_load_uint(cs, 32);
cs = __cs1;
let op = __r1?;
\end{verbatim}

and then \verb|__Slice_load_uint| lowers to the same \verb|LDUX|-based helper
family already used by current \verb|~load_uint| built-ins.

\subsection{Structs Lower as Products, Not Objects}
Structs do not imply heap allocation or object identity. Their canonical
meaning is just ``a named product type''. The lowering strategy is:

\begin{enumerate}
\item field names guide parsing, typing, and pattern matching;
\item a struct value is represented as a tuple in field order at observable
boundaries;
\item within one function, the backend may flatten fields into separate locals.
\end{enumerate}

This keeps structs zero-cost in the common case while still giving the user a
legible source-level model.

\subsection{Enums, Option, and Result Lower as Tagged Products}
General algebraic data types require an explicit discriminant. The canonical
runtime form is therefore:

\begin{verbatim}
(tag, payload)
\end{verbatim}

where:

\begin{itemize}
\item \verb|tag| is an integer discriminant;
\item \verb|payload| is a tuple of the selected variant fields, or null for a
zero-field variant.
\end{itemize}

Examples:

\begin{verbatim}
Option<T>::None        => (0, null)
Option<T>::Some(x)     => (1, x)
Result<T, E>::Ok(x)    => (0, x)
Result<T, E>::Err(e)   => (1, e)
\end{verbatim}

This representation is uniform, supports nesting, and matches ordinary
\verb|match| lowering. A backend may locally optimize some cases to a more
compact form, but the canonical semantic model should remain tagged.

\subsection{Pattern Matching Lowers to Tag Tests and Destructuring}
A \verb|match| expression lowers in four stages:

\begin{enumerate}
\item evaluate the scrutinee into a temporary;
\item extract or inspect its discriminant;
\item branch to the selected arm;
\item destructure the payload into arm-local bindings.
\end{enumerate}

For dense integer discriminants, the backend may generate a jump table. For
sparse discriminants, it may generate a branch tree. For message decoding, it
may first use a slice preload such as \verb|PLDUZ 32| or a full \verb|LDU 32|
to inspect the opcode before decoding the remainder of the body.

\subsection{The Question-Mark Operator Lowers to Early Return}
The \verb|?| operator should be understood as structured control flow, not as an
exception mechanism.

\begin{verbatim}
let msg = parse_msg(body)?;
\end{verbatim}

canonically lowers to:

\begin{enumerate}
\item evaluate \verb|parse_msg(body)| into a temporary \verb|Result|;
\item inspect the tag;
\item if the tag is \verb|Err|, package the error in the caller's return form
and jump to the function epilogue;
\item if the tag is \verb|Ok|, extract the payload and continue.
\end{enumerate}

This keeps all recoverable failure explicit in the generated control flow.

\subsection{Typed Messages Lower to Generated Codecs}
A declaration such as

\begin{verbatim}
pub message Transfer(op = 0x0f8a7ea5) {
    query_id: u64,
    amount: Coins,
    to: Address,
}
\end{verbatim}

should generate at least:

\begin{enumerate}
\item a stack-level runtime type equivalent to a struct;
\item an encoder that prepends \verb|0x0f8a7ea5| and then serializes fields;
\item a decoder that checks the prefix and then parses fields;
\item metadata allowing typed entrypoints and typed \verb|send| wrappers to
select the codec automatically.
\end{enumerate}

The encoder lowers to ordinary builder instructions such as \verb|NEWC|,
\verb|STU|/\verb|STUX|, address stores, amount stores, optional ref stores, and
\verb|ENDC|. The decoder lowers to \verb|CTOS|, \verb|LDU|/\verb|LDUX|,
address loads, amount loads, optional-value loads, and a final \verb|ENDS| when
the codec requires exact consumption.

\subsection{Typed Cell APIs Lower to Raw Cell Operations}
The typed cell APIs are wrappers over codecs and raw cell primitives:

\begin{itemize}
\item \verb|to_cell(x)| becomes ``encode \verb|x| into a builder, then
\verb|ENDC|'';
\item \verb|from_cell(c)| becomes ``\verb|CTOS|, decode fields, ensure final
slice invariants, return typed value'';
\item \verb|load_any<T>| becomes a statically resolved call to \verb|T|'s
decoder;
\item \verb|store_any<T>| becomes a statically resolved call to \verb|T|'s
encoder.
\end{itemize}

The important point is that \verb|Cell<T>| itself carries no runtime metadata.
All real work happens at codec boundaries.

\subsection{Typed Contexts Lower to ABI Wrappers}
A function such as

\begin{verbatim}
entry internal fn recv_internal(
    ctx: InternalContext<Transfer>,
) -> Result<(), WalletError> effect(read, write, send) { ... }
\end{verbatim}

should lower to two layers:

\begin{enumerate}
\item a generated low-level entry wrapper with the existing internal-message ABI;
\item a typed user function body that receives a constructed
\verb|InternalContext<Transfer>|.
\end{enumerate}

The wrapper:

\begin{enumerate}
\item reads the raw inbound context values from the existing ABI inputs;
\item extracts sender, transferred value, forward fee, and raw body slice;
\item decodes the raw body as \verb|Transfer|;
\item constructs the logical context tuple;
\item calls the typed user body;
\item converts the user body's result into the required low-level contract
behavior.
\end{enumerate}

This is protocol-aware lowering, but it is still only a wrapper around the real
VM inputs.

\subsection{Effect Sets Do Not Survive Code Generation}
An effect annotation such as

\begin{verbatim}
effect(read, write, send)
\end{verbatim}

should be enforced entirely by static analysis:

\begin{itemize}
\item a \verb|pure| function may not reach \verb|get_data|, \verb|set_data|,
\verb|SENDRAWMSG|, raw asm, or unsafe intrinsics;
\item a \verb|read|-only function may read context and storage but not write or
emit actions;
\item a \verb|send|-annotated function may use outbound-message primitives.
\end{itemize}

After validation, the annotation can be erased. The generated Fift code does not
need a runtime effect table.

\subsection{Traits and Generics Lower by Monomorphization}
Trait calls must never require a runtime interface object. The lowering model
is:

\begin{enumerate}
\item resolve the concrete implementation during type checking;
\item create a concrete instantiation of any generic function or method;
\item compile the resulting concrete function like any other function.
\end{enumerate}

For example, a call to:

\begin{verbatim}
send<Transfer>(dst, body, opts)
\end{verbatim}

should resolve the \verb|CellEncode| implementation for \verb|Transfer| at
compile time and produce a concrete body encoder call, not a reflective or
dynamic dispatch sequence.

\subsection{Typed Send Lowers to Ordinary Message Construction}
The high-level call

\begin{verbatim}
send(dst, body, opts)
\end{verbatim}

should lower to the following pipeline:

\begin{enumerate}
\item encode \verb|body| using the selected codec;
\item build the message envelope with the right destination, amount, and flags;
\item decide whether the body is inlined or stored by reference;
\item finalize the message cell;
\item emit \verb|SENDRAWMSG| with the chosen send mode.
\end{enumerate}

\verb|BounceMode| and \verb|BodyLayout| are therefore not new runtime entities.
They are source-level policy inputs to message construction.

\subsection{BodyLayout::Auto Must Be Deterministic}
\verb|BodyLayout::Auto| is useful only if all compliant toolchains make the same
choice for the same source and codec. The canonical rule should therefore be:

\begin{enumerate}
\item compute the encoded body size exactly when it is statically known;
\item otherwise encode into a temporary builder/cell first;
\item choose inline vs ref according to a fully specified fit predicate;
\item emit the corresponding branch or specialized fast path.
\end{enumerate}

The exact fit predicate belongs to the language+stdlib contract and must be
stable across compilers.

\section{Worked Lowering Example}
\subsection{Source Fragment}
\begin{verbatim}
pub message Transfer(op = 0x0f8a7ea5) {
    seqno: u32,
    to: Address,
    amount: Coins,
}

entry internal fn recv_internal(
    ctx: InternalContext<Transfer>,
) -> Result<(), WalletError> effect(read, write, send) {
    let mut state = load_state()?;
    let msg = ctx.body;

    if msg.seqno != state.seqno {
        return Err(WalletError::WrongSeqno);
    }

    send(
        msg.to,
        (),
        SendOptions {
            value: msg.amount,
            bounce: BounceMode::Body256,
            body_layout: BodyLayout::Auto,
            mode: SendMode::PayFeesSeparately,
        }
    )?;

    state.seqno += 1;
    save_state(state)?;
    Ok(())
}
\end{verbatim}

\subsection{Conceptual Lowering}
The canonical lowering is:

\begin{enumerate}
\item generate \verb|__Transfer_encode| and \verb|__Transfer_decode|;
\item generate a low-level internal-entry wrapper \verb|__entry_recv_internal|;
\item keep the typed user body as a concrete selector function;
\item lower each \verb|?| into explicit branch-and-return control flow;
\item lower \verb|send(... )| to ordinary message building plus
\verb|SENDRAWMSG|.
\end{enumerate}

\subsection{Generated Wrapper Shape}
\begin{verbatim}
__entry_recv_internal(low_level_abi_args):
    raw_ctx   = parse_low_level_internal_context(low_level_abi_args)
    msg_res   = __Transfer_decode(raw_ctx.body_slice)
    if msg_res is Err(e):
        return lower_entry_parse_failure(e)
    typed_ctx = InternalContext {
        sender: raw_ctx.sender,
        value: raw_ctx.value,
        forward_fee: raw_ctx.forward_fee,
        body: msg_res.ok,
    }
    return __recv_internal_impl(typed_ctx)
\end{verbatim}

\subsection{User Body Shape After Desugaring}
\begin{verbatim}
__recv_internal_impl(ctx):
    __r0 = load_state()
    if __r0 is Err(e): return Err(WalletError::Parse(e))
    state = __r0.ok
    msg = ctx.body

    if msg.seqno != state.seqno:
        return Err(WalletError::WrongSeqno)

    __r1 = send(msg.to, (), opts_from(msg.amount))
    if __r1 is Err(e): return Err(WalletError::Send(e))

    state.seqno = state.seqno + 1

    __r2 = save_state(state)
    if __r2 is Err(e): return Err(map_cell_error(e))

    return Ok(())
\end{verbatim}

\subsection{Fift/TVM Shape}
At the final backend stage, the remaining building blocks are ordinary TVM
operations:

\begin{itemize}
\item state load/save via the current persistent-data path;
\item integer comparison for the sequence number check;
\item builder construction and \verb|ENDC| for any encoded cells;
\item \verb|SENDRAWMSG| for the outbound transfer;
\item selector-based function calls for helper routines;
\item conditional branches for all \verb|Result| tests.
\end{itemize}

In other words, the high-level source has been reduced to the same machine model
already used by present FunC and hand-written Fift.

\section{Required Canonical Decisions Before Implementation}
The proposal is lowerable, but several details must be fixed precisely before a
compiler implementation can claim compatibility:

\begin{enumerate}
\item the exact canonical stack representation of \verb|enum|,
\verb|Option<T>|, and \verb|Result<T, E>|;
\item the exact derived wire format for serialized enums and results;
\item whether \verb|Bytes| is canonically a byte-aligned slice, a typed cell,
or a stdlib wrapper with a fixed ABI;
\item the exact deterministic rule behind \verb|BodyLayout::Auto|;
\item the mapping from typed error results in entry wrappers to legacy numeric
exit behavior where required;
\item the exact getter/export selector policy for edition \verb|0.5.0|.
\end{enumerate}

These are not objections to the design. They are the finite set of ABI-level
decisions required to keep different compilers interoperable.

\section{Conclusion}
FunC \verb|0.5.0| can remain fully lowerable to Fift/TVM if it preserves one
discipline: \emph{push complexity into the front end, not into the runtime}.

Under that discipline:

\begin{itemize}
\item modules, visibility, effects, and mutable-place checks erase entirely;
\item structs, enums, messages, and results become fixed tuple/cell layouts;
\item typed entrypoints and typed send wrappers become generated adapters over
existing protocol conventions;
\item traits and generics remain static-dispatch only;
\item the final emitted operations are still the familiar TVM primitives for
control flow, parsing, building, state access, and \verb|SENDRAWMSG|.
\end{itemize}

This is exactly the kind of language upgrade that keeps TVM/TOS honest:
better source ergonomics, stronger static guarantees, and no semantic break
from the underlying machine.

\end{document}
